\input{../../../assets/programming-editorial.tex}
\renewcommand{\studiumcoverphoto}{../../../assets/covers/vol06.jpg}
\urldef{\studiumsourceurl}\url{https://commons.wikimedia.org/wiki/File:Linux_Workstation.jpg}
\urldef{\studiumlicenseurl}\url{http://creativecommons.org/publicdomain/zero/1.0/deed.en}
\newcommand{\studiumcustomtitle}{
\studiumcover{JVM, .NET, Go\\ y desarrollo móvil}{Máquinas virtuales, ecosistemas y plataformas}{VOLUMEN 6}
\chapter*{Créditos y alcance de esta revisión}
Daniel Miguel Tejedor. Biblioteca de Programación, Studium. 2026.

Esta copia conserva los 75 capítulos de la colección original. La revisión de fuentes, cobertura y soluciones sigue abierta. Un cálculo aritmético reproducido no verifica un programa. Esta edición no cuenta con aprobación académica independiente.

\noindent\textbf{Fotografía de portada.} K.Sreeram Wiki / Wikimedia Commons / CC0. Linux Workstation.jpg.\par\studiumsourceurl

\noindent\textbf{Licencia y transformaciones.} CC0.\par\studiumlicenseurl

Original bytes unchanged. Proportional placement and crop in shared LaTeX cover. La imagen y su recorte conservan la licencia indicada.

\noindent\textbf{Cómo estudiar.} Escribe los ejemplos en archivos, ejecuta las pruebas indicadas y compara la salida observada con la esperada. Consulta las soluciones después de intentarlo. Los bloques Terminal contienen comandos; Output muestra resultados y Diagnostic contiene errores.\par\clearpage
}
